\section{Points of schemes}
\label{section-points}

\noindent
Given a scheme $X$ we can define a functor
$$
h_X : \Sch^{opp}
\longrightarrow
\textit{Sets}, \quad
T \longmapsto \Mor(T, X).
$$
See Categories, Example \ref{categories-example-hom-functor}.
This is called the {\it functor of points of $X$}.
A fun part of scheme theory is to find descriptions of
the internal geometry of $X$ in terms of this functor $h_X$.
In this section we find a simple way to describe
points of $X$.

\medskip\noindent
Let $X$ be a scheme. Let $R$ be a local ring with maximal ideal
$\mathfrak m \subset R$. Suppose that $f : \Spec(R) \to X$
is a morphism of schemes. Let $x \in X$ be the image of the closed point
$\mathfrak m \in \Spec(R)$. Then we obtain a local homomorphism
of local rings
$$
f^\sharp :
\mathcal{O}_{X, x}
\longrightarrow
\mathcal{O}_{\Spec(R), \mathfrak m} = R.
$$

\begin{lemma}
\label{lemma-morphism-from-spec-local-ring}
Let $X$ be a scheme. Let $R$ be a local ring.
The construction above gives a bijective correspondence
between morphisms $\Spec(R) \to X$ and pairs
$(x, \varphi)$ consisting of a point $x \in X$ and
a local homomorphism of local rings $\varphi : \mathcal{O}_{X, x} \to R$.
\end{lemma}

\begin{proof}
Let $A$ be a ring. For any ring homomorphism $\psi : A \to R$
there exists a unique prime ideal $\mathfrak p \subset A$
and a factorization $A \to A_{\mathfrak p} \to R$ where the
last map is a local homomorphism of local rings. Namely,
$\mathfrak p = \psi^{-1}(\mathfrak m)$. Via
Lemma \ref{lemma-morphism-into-affine}
this proves that the lemma holds if $X$ is an affine scheme.

\medskip\noindent
Let $X$ be a general scheme. Any $x \in X$ is contained in
an open affine $U \subset X$. By the affine case we conclude that every pair
$(x, \varphi)$ occurs as the end product of the construction
above the lemma.

\medskip\noindent
To finish the proof it suffices to show that any morphism
$f : \Spec(R) \to X$ has image contained in any affine
open containing the image $x$ of the closed
point of $\Spec(R)$. In fact, let $x \in V \subset X$
be any open neighbourhood containing $x$. Then
$f^{-1}(V) \subset \Spec(R)$ is an open containing
the unique closed point and hence equal to $\Spec(R)$.
\end{proof}

\noindent
As a special case of the lemma above we obtain for every
point $x$ of a scheme $X$ a canonical morphism
\begin{equation}
\label{equation-canonical-morphism}
\Spec(\mathcal{O}_{X, x}) \longrightarrow X
\end{equation}
corresponding to the identity map on the local ring of $X$ at $x$.
We may reformulate the lemma above as saying that for any
morphism $f : \Spec(R) \to X$ there exists a unique point
$x \in X$ such that $f$ factors as
$\Spec(R) \to \Spec(\mathcal{O}_{X, x}) \to X$
where the first map comes from a local homomorphism
$\mathcal{O}_{X, x} \to R$.

\medskip\noindent
In case we have a morphism of schemes $f : X \to S$,
and a point $x$ mapping to a point $s \in S$ we obtain
a commutative diagram
$$
\xymatrix{
\Spec(\mathcal{O}_{X, x}) \ar[r] \ar[d] & X \ar[d] \\
\Spec(\mathcal{O}_{S, s}) \ar[r] & S
}
$$
where the left vertical map corresponds to the local ring map
$f^\sharp_x : \mathcal{O}_{S, s} \to \mathcal{O}_{X, x}$.

\begin{lemma}
\label{lemma-specialize-points}
Let $X$ be a scheme.
Let $x, x' \in X$ be points of $X$.
Then $x' \in X$ is a generalization of $x$ if and only if
$x'$ is in the image of the canonical morphism
$\Spec(\mathcal{O}_{X, x}) \to X$.
\end{lemma}

\begin{proof}
A continuous map preserves the relation of specialization/generalization.
Since every point of $\Spec(\mathcal{O}_{X, x})$ is a
generalization of the closed point we see every point in the image
of $\Spec(\mathcal{O}_{X, x}) \to X$ is a generalization of $x$.
Conversely, suppose that $x'$ is a generalization of $x$.
Choose an affine open neighbourhood $U = \Spec(R)$ of
$x$. Then $x' \in U$. Say $\mathfrak p \subset R$ and
$\mathfrak p' \subset R$ are the primes corresponding
to $x$ and $x'$. Since $x'$ is a generalization of $x$
we see that $\mathfrak p' \subset \mathfrak p$. This means
that $\mathfrak p'$ is in the image of the morphism
$\Spec(\mathcal{O}_{X, x}) = \Spec(R_{\mathfrak p})
\to \Spec(R) = U \subset X$ as desired.
\end{proof}

\noindent
Now, let us discuss morphisms from spectra of fields.
Let $(R, \mathfrak m, \kappa)$ be a local ring
with maximal ideal $\mathfrak m$ and residue field $\kappa$.
Let $K$ be a field. A local homomorphism $R \to K$ by definition
factors as $R \to \kappa \to K$, i.e., is the same thing as
a morphism $\kappa \to K$. Thus we see that morphisms
$$
\Spec(K) \longrightarrow X
$$
correspond to pairs $(x, \kappa(x) \to K)$. We may define
a preorder on morphisms of spectra of fields to $X$
by saying that $\Spec(K) \to X$ dominates
$\Spec(L) \to X$ if $\Spec(K) \to X$
factors through $\Spec(L) \to X$. This suggests
the following notion: Let us temporarily
say that two morphisms $p : \Spec(K) \to X$ and
$q : \Spec(L) \to X$ are {\it equivalent} if there exists
a third field $\Omega$ and a commutative diagram
$$
\xymatrix{
\Spec(\Omega) \ar[r] \ar[d] &
\Spec(L) \ar[d]^q \\
\Spec(K) \ar[r]^p &
X
}
$$
Of course this immediately implies that the unique points of
all three of the schemes $\Spec(K)$,
$\Spec(L)$, and $\Spec(\Omega)$
map to the same $x \in X$. Thus a diagram (by the remarks above)
corresponds to a point $x \in X$ and a commutative diagram
$$
\xymatrix{
\Omega &
L \ar[l] \\
K \ar[u] &
\kappa(x) \ar[l] \ar[u]
}
$$
of fields. This defines an equivalence relation, because given
any set of field extensions $K_i/\kappa$ there exists
some field extension $\Omega/\kappa$ such that all
the field extensions $K_i$ are contained in the extension $\Omega$.

\begin{lemma}
\label{lemma-characterize-points}
Let $X$ be a scheme. Points of $X$ correspond bijectively
to equivalence classes of morphisms from spectra of
fields into $X$. Moreover, each equivalence class contains
a (unique up to unique isomorphism) smallest element
$\Spec(\kappa(x)) \to X$.
\end{lemma}

\begin{proof}
Follows from the discussion above.
\end{proof}

\noindent
Of course the morphisms $\Spec(\kappa(x)) \to X$
factor through the canonical morphisms
$\Spec(\mathcal{O}_{X, x}) \to X$.
And the content of Lemma \ref{lemma-specialize-points} is in
this setting that the morphism $\Spec(\kappa(x')) \to X$
factors as
$\Spec(\kappa(x')) \to \Spec(\mathcal{O}_{X, x}) \to X$
whenever $x'$ is a generalization of $x$.
In case we have a morphism of schemes $f : X \to S$,
and a point $x$ mapping to a point $s \in S$ we obtain
a commutative diagram
$$
\xymatrix{
\Spec(\kappa(x)) \ar[r] \ar[d] &
\Spec(\mathcal{O}_{X, x}) \ar[r] \ar[d] &
X \ar[d] \\
\Spec(\kappa(s)) \ar[r] &
\Spec(\mathcal{O}_{S, s}) \ar[r] &
S.
}
$$










